\documentclass[9pt]{beamer-control}
\usepackage{beamer-control-prac}
\begin{document}
\CONCEPT[7]{Weeks 7--9: Complex system --- Ping pong rig}

\begin{frame}
\frametitle{Introduction}
In this practical project, we will implement a PID controller on a new piece of hardware, the ping pong rig. This plant emphasises the versatility of the PID control methodology. 

\vfill

This practical will consist of the following parts:
\begin{itemize}
\item Setup
\item Calibration
\item Characterisation
\item Tuning
\item Report
\end{itemize}
\end{frame}

\SUBCONCEPT{Setup}

\begin{frame}{Ping pong rig hardware}
\begin{columns}[T]
\begin{column}{0.7\textwidth}
\vspace{1cm}
\begin{itemize}
\item The ping pong rig has one input and one output
\item The input is a voltage supplied to a fan mounted on the base of the system to create an airflow to push the ping pong ball up the vertical cylinder
\item The output is a voltage reading from an ultrasonic proximity sensor placed at the bottom of the tube facing upwards, which will need to be converted into units of distance for control purposes
\item The ultrasonic sensor is less accurate and more noisy than the encoder used in previous practicals, making it more of a challenge to adequately control
\end{itemize}
\end{column}
\begin{column}{0.3\textwidth}
		\begin{figure}
		\centering
		\includegraphics[width=2.5cm]{project_pingpongrig}
	\end{figure}
\end{column}
\end{columns}
\end{frame}


\begin{frame}{Setting up hardware}
Get into groups of 2-3 and follow these steps
\begin{enumerate}
	\item Connect the rig to the power supply
	\item Connect the ping pong rig output labelled Distance to channel AI 0 on the board
	\item Connect the ping pong rig input labelled Fan Input to channel AO 0 in the board
	\item Ensure the knob gain is rotated all the way to the right (minimum gain)
\end{enumerate}
\end{frame}

\begin{frame}{Setting up software}
	Open Simulink and create a blank model. In this practical you are going to create your own model for the ping pong rig. There are three blocks that initiate and receive/send signals from the physcial rig:
	\begin{itemize}
		\item HIL Intialize: Hardware Initialization
		\item HIL Read Analog: Receive analogue signal (voltage)
		\item HIL Write Analog: Send analogue signal (voltage)
	\end{itemize}
	Find these three blocks from the LIbrary Browser and drag them into your model. Double click on the HIL Read/Write blocks and ensure channel 0 is selected. Double click on HIL Initialize and choose the PCIE 6321 board under Board Type Section. 
\end{frame}

\begin{frame}{Setting up software}
	Select the three blocks then right click and select create a subsystem. Now this is your ping
pong plant which has an input and an output, you can label them as shown below. \\
\vspace{0.5cm}

The \textit{To Workspace} block exports the position data to MATLAB. Set up the \textit{Step} block with a
step time of 0 and a final value of between 1 and 10 volts and run. You may need to
increase the input voltage to increase the upward force provided from the fan.\\

\begin{figure}
	\centering
	\includegraphics[width=10cm]{project_subsystem}
\end{figure}
\end{frame}


\begin{frame}{Running the system}
	To run on Hardware go to Modelling > Model Setting > Code Generation > Browse > Choose quar.win64.tlc, now the Hardware tab should appear.\\
	\vspace{0.5cm}
	
	The To Workspace block may work incorrectly on your model, where data points are lost
	when the model is running. To prevent this, go to Modelling > Model Setting > Solver >
	Solver details > Change fixed-step size to auto (or 0.01). \\
	\vspace{0.5cm}
	Try running the system to get the ping pong ball to lift into the tube by adding a step input to the system. You may have to play around with the voltage input to get the ball to lift. \\
	\vspace{0.5cm}
	The system will inherit the final voltage between runs and so it might be beneficial to change the Stop Time to 10 seconds, Step Time to 9 seconds, Initial Value to your desired input voltage, and Final Value to a low voltage that doesn't lift the ball (such as 2V). 
	
	
\end{frame}

\SUBCONCEPT{Calibration}

\begin{frame}{Voltage to distance conversion}
	After getting the ping pong ball to lift, you may notice that the output is not in units of distance. You will need to calibrate the sensor (i.e. convert volts to a unit of distance).
\begin{itemize}
\item Measure the \textit{output} sensor voltage (not the input voltage!) with its corresponding distance by using the 10cm markings on the side of the tube
\item Use two data points to construct a linear relationship to correlate the output sensor voltage with distance
\item Use the Gain and Constant blocks to implement this conversion
\end{itemize}
\begin{figure}
	\centering
	\includegraphics[width=10cm]{project_conversion}
\end{figure}
\end{frame}


\SUBCONCEPT{Characterisation}

\begin{frame}{Step response}
Now that the system has been calibrated, we may characterise the response of the plant.
	\begin{itemize}
		\item Run a step response on the rig
		\item Plot the data in MATLAB using the To Workspace block and attempt to fit a straight line to find $a$ and $\tau$
	\end{itemize}
You will notice that the sensor is heavily discretised, making analysis difficult.
\end{frame}

\begin{frame}{Processing the output}
Due to the discretisation, the derivative may give large spikes which would limit the
effectiveness of a controller. Therefore, as with previous practicals, we must filter the output
to smooth the signal without introducing too much time lag. Filter the displacement using the
same method in previous practicals to identify an appropriate cut off frequency. An example
of a filtered signal is shown below. Clearly, the signal is smoother, does not have too much
lag associated and may be suitable for controlling purposes. 

\begin{figure}
	\centering
	\includegraphics[width=6cm]{project_filter}
\end{figure}
\end{frame}


\SUBCONCEPT{Tuning}

\begin{frame}{Inital tuning}
	Now that the system has been appropriately setup and calibrated, and the output has been conditioned using a lowpass filter and characterised using an ITD model we may start tuning.
	\begin{itemize}
		\item Using the values for $a$ and $\tau$, find PID gains
		\item Implement a closed loop PID controller in Simulink
	\end{itemize}
\begin{figure}
	\centering
	\includegraphics[width=11cm]{project_pid}
\end{figure}
\end{frame}


\begin{frame}{Manual tuning}
With these initial gains:
\begin{itemize}
	\item Does the controller stabilise the system?
	\item Does it achieve setpoint following?
	\item How does the controller perform in general?
\end{itemize}
Your initial controller was likely not performing very well. Even if it was, manual tuning is necessary to increase the performance of it.
\end{frame}

\begin{frame}{Performance quantification}
	Once you have successfully controlled the ping pong rig, try varying some of the system parameters, adding disturbance into the system, or trying set point following. \\
	\vspace{0.5cm}
	This investigation will indicate the robustness of your controller and how well your system operates under different conditions.\\
	\vspace{0.5cm}
	This part of the project is purposefully open-ended, what do \textit{you} think are important experiments to test the performance of the ping pong rig controller?
\end{frame}


\SUBCONCEPT{Report}

\begin{frame}{Project report}
	After designing and testing your controller you will describe your process in a professional report. While the structure is up to you, a few suggestions for inclusion are:
	\begin{itemize}
		\item A description of the hardware and software
		\item A description of the calibration process
		\item A discussion on open loop response characteristics such as rise time, overshoot etc.
		\item An explanation of the tuning process such as what tuning law was used, your method of manual tuning, etc.
		\item An in-depth discussion of the performance of the controller
		\item Diagrams of Simulink models and MATLAB code placed in an appendix
	\end{itemize}
	Please use the To Workspace block for high quality plots in MATLAB rather than screenshots of Simulink.
\end{frame}


\SUMMARYFRAME

\end{document}
